% =============================================================================
%  Preambulo: pacotes, estilos e configuracoes globais.
%  As variaveis \mytitle / \myauthors / \mydate sao definidas em report.tex
%  ANTES deste arquivo ser carregado.
% =============================================================================

% MASA Formal Report Template frontmatter, conforming to AE405 Standards
% Created by Theo Rulko, 25/12/2021
% Contact trulko@umich.edu with questions

% ------------------------------------------
% Import primary packages
% O \documentclass fica em report.tex, e nao aqui: e o unico jeito de o
% Overleaf identificar sozinho o arquivo mestre na importacao do .zip.
% Ver tools/LEIA-ME-OVERLEAF.txt.
\usepackage{datetime}
\usepackage[T1]{fontenc}
\usepackage{graphicx,algorithmic,algorithm}
% xcolor (superconjunto de color): o tcolorbox o carrega mais adiante, e
% carrega-lo aqui evita 'Incompatible color definition' em MASA-Blue.
\usepackage{xcolor}
\usepackage{listings,url}
\usepackage{pdfpages}
\usepackage{wallpaper}
\usepackage{titlesec}
\usepackage{enumitem}
\usepackage{changepage}
\usepackage{pdfpages}
\usepackage{lscape}
\usepackage{amsmath}
\usepackage{montserrat}
\usepackage{lipsum}

% ------------------------------------------
% ToC, AIAA-style Bibliography
\usepackage[immediate]{silence}
\WarningFilter{latex}{You have requested package}
\usepackage[tocgraduated]{assets/tocstyle}
\usetocstyle{standard}
\RequirePackage[sort&compress,numbers]{natbib}
\bibliographystyle{new-aiaa}
\WarningFilter[temp]{latex}{Command}
    \usepackage{sectsty}
\DeactivateWarningFilters[temp]
\usepackage[nottoc]{tocbibind}

% ------------------------------------------
% Figures & tables
\usepackage{float}
\usepackage[small,bf]{caption}
\usepackage{tabularx}
\usepackage{multirow}
\usepackage{placeins}
\usepackage{makecell}
\usepackage[figuresleft]{rotating}

% ------------------------------------------
% Custom layout parameters
\usepackage[margin=0.75in]{geometry}
\definecolor{MASA-Blue}{RGB}{0,39,76}
\renewcommand\thesection{\Roman{section}}
\renewcommand{\bibsection}{\section{References}}
\newcommand{\boldheader}[1]{\vspace{5mm}\textbf{#1}\textbf{:}\newline}
\sectionfont{\fontsize{19}{20}\color{MASA-Blue}\fontseries{b}\fontfamily{Montserrat-TOsF}\selectfont}
\subsectionfont{\fontsize{16}{18}\color{MASA-Blue}\fontfamily{Montserrat-TOsF}\selectfont}
\subsubsectionfont{\fontseries{sb}\fontshape{n}\color{MASA-Blue}\fontfamily{Montserrat-TOsF}\selectfont}
\titlespacing{\subsection}{0pt}{100mm}{0pt}
\setlength{\parindent}{0in}
\setlength{\parskip}{1em}

% ------------------------------------------
% Letterhead
\newcommand\MASALetterhead[3]{
    \hfuzz=60pt
    \thispagestyle{empty}
    \setlength{\parskip}{0in}
    % (For full style of letterhead footer, use Footer1)
    \LLCornerWallPaper{1}{assets/Footer2.pdf}
    \begin{flushleft}
        \includegraphics[width=4.5cm]{assets/Worm_vector.pdf}
    \end{flushleft}
    \begin{flushright}
        \begin{tabularx}{0.5\textwidth}{l @{}}#1\end{tabularx}
    \end{flushright}
    \begin{flushleft}
        \par #2 \\ \vspace{2mm} \mydate
    \end{flushleft}
    Subject: #3
    \setlength{\parskip}{1em}
}
\newcommand\signature[2]{
    \begin{minipage}{5cm}
        \noindent\vspace{4.5mm}\par
        \noindent\vspace{-1mm}\rule{#1cm}{0.5pt}\par
        \noindent#2
    \end{minipage}
}

% ------------------------------------------
% Title Page
\newcommand\makeTitlePage[1]{
    \newpage
    \ClearWallPaper
    \thispagestyle{empty}
    \setlength{\parskip}{0in}
    % (For TSM logo footer, use Footer3)
    \LLCornerWallPaper{1}{assets/Footer4.pdf}
    \vspace*{\fill}
    \begin{adjustwidth}{1.5cm}{1.5cm} 
    \begin{center}
        % \includegraphics[width=7cm]{assets/Worm_vector.pdf}\\
        % \vspace{20mm}
        \fontsize{28}{30}\fontseries{b}\fontshape{n}\color{MASA-Blue}\fontfamily{Montserrat-TOsF}\selectfont{\mytitle}\\
        \vspace{15mm}
        \fontsize{14}{16}\fontseries{sb}\fontshape{n}\color{MASA-Blue}\fontfamily{Montserrat-TOsF}\selectfont{\mydate\\\vspace{7.5mm} Prepared by:}\\
        \fontsize{14}{16}\fontseries{m}\fontshape{n}\color{MASA-Blue}\fontfamily{Montserrat-TOsF}\selectfont{\myauthors}\\
        \vspace{7.5mm}
        \fontsize{14}{16}\fontseries{sb}\fontshape{n}\color{MASA-Blue}\fontfamily{Montserrat-TOsF}\selectfont{Prepared for:}\\
        \fontsize{14}{16}\fontseries{m}\fontshape{n}\color{MASA-Blue}\fontfamily{Montserrat-TOsF}\selectfont{#1}
        \vspace{50mm}
    \end{center}
    \end{adjustwidth}
    \vspace*{\fill}
    \setlength{\parskip}{1em}
}

% ------------------------------------------
% Cover Page
\newcommand\makeCoverPage{
    \newpage
    \ClearWallPaper
    \thispagestyle{empty}
    \setlength{\parskip}{0in}
    % (For TSM logo footer, use Footer3)
    \LLCornerWallPaper{1}{assets/Footer4.pdf}
    \vspace*{\fill}
    \begin{adjustwidth}{1.5cm}{1.5cm} 
    \begin{center}
        \includegraphics[width=9cm]{assets/Worm_vector.pdf}\\
        \vspace{25mm}
        \fontsize{33}{35}\fontseries{b}\fontshape{n}\color{MASA-Blue}\fontfamily{Montserrat-TOsF}\selectfont{\mytitle}
        \vspace{45mm}
    \end{center}
    \end{adjustwidth}
    \vspace*{\fill}
    \setlength{\parskip}{1em}
}

% ------------------------------------------
% Acknowledgements
\newcommand\MASAacknowledgements[1]{
    \newpage
    \topskip0pt
    \addcontentsline{toc}{section}{Acknowledgements}
    \vspace*{\fill}
    \begin{adjustwidth}{1.5cm}{1.5cm} 
    \begin{center}
    \textbf{\Large\color{MASA-Blue}\fontfamily{Montserrat-TOsF}\selectfont Acknowledgements}\\
    \vspace{5mm}
    \end{center}
    #1
    \end{adjustwidth}
    \vspace*{\fill}
}

% ------------------------------------------
% Format for matlab code listing
\usepackage{listings}
\definecolor{commentgreen}{RGB}{14,120,0}
\lstset{
    numbers=left, 
    basicstyle=\scriptsize, 
    commentstyle=\color{commentgreen},
    keywordstyle=\bf\color{blue},
    language=matlab, 
    frame=leftline,
    rulesepcolor=\color{black},
    flexiblecolumns=true,
    extendedchars=false,
    keepspaces=true
}

% ------------------------------------------
% Custom Math Shorthand (Credit: kfid@umich.edu)
\newcommand{\pp}[2]{\dfrac{\partial #1}{\partial #2}}
\newcommand{\ppt}[2]{\dfrac{\partial^2 #1}{\partial #2^2}}
\newcommand{\DD}[2]{\dfrac{D #1}{D #2}}
\newcommand{\dd}[2]{\dfrac{d #1}{d #2}}
\newcommand{\mbb}[1]{\mathbb{#1}}
\newcommand{\mbf}[1]{\mathbf{#1}}
\newcommand{\mrm}[1]{\mathrm{#1}}
\newcommand{\mcal}[1]{\mathcal{#1}}
\newcommand{\BigO}{\mathcal{O}}
\newcommand{\norm}[2]{\left\lVert{#1}\right\lVert_{#2}}
\newcommand{\sbf}[1]{\boldsymbol{#1}}
\def\etal{{\it et al~}}
\newcommand{\be}{\begin{eqnarray}}
\newcommand{\ee}{\end{eqnarray}}
\newcommand{\ben}{\begin{eqnarray*}}
\newcommand{\een}{\end{eqnarray*}}
\newcommand{\dx}{d\mbf{x}}
\newcommand{\Dt}{\Delta t}
\newcommand{\ih}{\hat{i}}
\newcommand{\jh}{\hat{j}}
\newcommand{\kh}{\hat{k}}
\newcommand{\xh}{\hat{x}}
\newcommand{\yh}{\hat{y}}
\newcommand{\zh}{\hat{z}}
\newcommand{\Uinf}{U_{\infty}}
\newcommand{\Vinf}{V_{\infty}}
\newcommand{\Minf}{M_{\infty}}
\newcommand{\bVinf}{\mbf{V}_{\infty}}
\newcommand{\scriptth}{\scriptsize \textnormal{th}}
\newcommand{\half}{\frac{1}{2}}
\newcommand{\CFL}{{\mbox{CFL}}}
\newcommand{\tn}[1]{\textnormal{#1}}
\newcommand{\mbfh}[1]{\widehat{\mathbf{#1}}}
\newcommand{\m}{\textrm{m}}
\newcommand{\s}{\textrm{s}}
\newcommand{\kg}{\textrm{kg}}
\newcommand{\km}{\textrm{km}}
\newcommand{\N}{\textrm{N}}
\newcommand{\Pa}{\textrm{Pa}}
\newcommand{\K}{\textrm{K}}
\newcommand{\C}{\textrm{C}}
\newcommand{\mbfv}[1]{\vec{\mathbf{#1}}}
\newcommand{\comment}[1]{}

% Paleta usada pelo estilo de listagem `bt`.
\definecolor{codekw}{RGB}{0,39,76}        % palavras-chave  (MASA-Blue)
\definecolor{codecomment}{RGB}{14,120,0}  % comentarios e strings
\definecolor{codegray}{RGB}{245,245,245}  % fundo (opcional)

% Sub-figuras (\subcaption dentro de minipages).
\usepackage{subcaption}

% =============================================================================
%  Linguagens que o pacote listings nao conhece de fabrica.
%  Sem estas definicoes, `language=Meson` e `language=GAS` abortam a
%  compilacao com "Couldn't load requested language".
% =============================================================================

% Meson -- a DSL do sistema de build usado por todo o repositorio. O listings
% nao a conhece; sem esta definicao, `language=Meson` aborta a compilacao com
% "Couldn't load requested language". As palavras destacadas sao as funcoes
% embutidas do Meson e as palavras-chave da linguagem, de modo que o nome de um
% alvo do repositorio apareca em texto normal -- a distincao que mais importa
% ao ler um meson.build.
\lstdefinelanguage{Meson}{
    morekeywords={project,executable,library,static_library,shared_library,
                  dependency,declare_dependency,subdir,subproject,
                  include_directories,files,find_program,import,
                  configure_file,custom_target,run_command,run_target,
                  test,benchmark,install_data,install_headers,install_subdir,
                  get_option,add_project_arguments,add_global_arguments,
                  message,warning,error,assert,summary,
                  if,elif,else,endif,foreach,endforeach,and,or,not,in,
                  true,false,continue,break},
    sensitive=true,
    morecomment=[l]{\#},
    morestring=[b]',
    morestring=[b]",
    alsoletter={_},
}

% Assembly GNU/AT&T (usado nos capitulos de baixo nivel). O dialeto embutido do
% listings e o da Intel/MASM; aqui interessam as diretivas do GAS e o
% mnemonico com sufixo de tamanho (movq, addl), que o dialeto padrao nao cobre.
\lstdefinelanguage{GAS}{
    morekeywords={mov,movq,movl,movb,movw,lea,leaq,add,addq,addl,sub,subq,
                  imul,imulq,idiv,cmp,cmpq,cmpl,test,testq,jmp,je,jne,jl,jle,
                  jg,jge,ja,jb,call,ret,push,pushq,pop,popq,xor,xorq,and,or,
                  shl,shr,sar,cmove,cmovg,cmovl,rep,movsb,syscall,cpuid,rdtsc,
                  nop,endbr64,vaddps,vmulps,vmovups},
    morekeywords=[2]{.text,.data,.bss,.globl,.global,.type,.size,.section,
                     .align,.quad,.long,.byte,.asciz,.string,.intel_syntax,
                     .att_syntax,.cfi_startproc,.cfi_endproc},
    sensitive=true,
    morecomment=[l]{\#},
    morecomment=[l]{//},
    morestring=[b]",
    alsoletter={.\_},
}

% -----------------------------------------------------------------------------
%  Caracteres nao-ASCII dentro de lstlisting
%
%  O listings processa o codigo byte a byte. Num arquivo UTF-8 uma letra
%  acentuada ocupa dois bytes, e sem tratamento eles sao compostos fora de
%  ordem: "integracao" acentuado sai como "cãintegrao". A tabela `literate`
%  abaixo mapeia cada caractere para o comando LaTeX equivalente, contando
%  como uma unica coluna para nao quebrar o alinhamento do codigo.
%  Ao usar um caractere novo dentro de um bloco de codigo, acrescente-o aqui.
% -----------------------------------------------------------------------------
\lstdefinestyle{btchars}{
    extendedchars = true,
    literate      =
      % acentuacao portuguesa
      {á}{{\'a}}1 {à}{{\`a}}1 {â}{{\^a}}1 {ã}{{\~a}}1 {ä}{{\"a}}1
      {é}{{\'e}}1 {è}{{\`e}}1 {ê}{{\^e}}1 {ë}{{\"e}}1
      {í}{{\'i}}1 {ì}{{\`i}}1 {î}{{\^i}}1 {ï}{{\"i}}1
      {ó}{{\'o}}1 {ò}{{\`o}}1 {ô}{{\^o}}1 {õ}{{\~o}}1 {ö}{{\"o}}1
      {ú}{{\'u}}1 {ù}{{\`u}}1 {û}{{\^u}}1 {ü}{{\"u}}1
      {ç}{{\c c}}1 {ñ}{{\~n}}1
      {Á}{{\'A}}1 {À}{{\`A}}1 {Â}{{\^A}}1 {Ã}{{\~A}}1 {Ä}{{\"A}}1
      {É}{{\'E}}1 {Ê}{{\^E}}1 {Í}{{\'I}}1
      {Ó}{{\'O}}1 {Ô}{{\^O}}1 {Õ}{{\~O}}1
      {Ú}{{\'U}}1 {Ü}{{\"U}}1 {Ç}{{\c C}}1
      {º}{{\textordmasculine}}1 {ª}{{\textordfeminine}}1
      % simbolos matematicos e gregos
      {→}{{$\rightarrow$}}1 {←}{{$\leftarrow$}}1
      {≈}{{$\approx$}}1 {≤}{{$\leq$}}1 {≥}{{$\geq$}}1 {≠}{{$\neq$}}1
      {×}{{$\times$}}1 {−}{{$-$}}1 {·}{{$\cdot$}}1 {…}{{$\dots$}}1
      {°}{{$^\circ$}}1 {²}{{$^2$}}1 {³}{{$^3$}}1
      {µ}{{$\mu$}}1 {§}{{\S}}1
      {α}{{$\alpha$}}1 {β}{{$\beta$}}1 {Δ}{{$\Delta$}}1 {θ}{{$\theta$}}1
      {λ}{{$\lambda$}}1 {μ}{{$\mu$}}1 {π}{{$\pi$}}1 {σ}{{$\sigma$}}1
      {τ}{{$\tau$}}1 {φ}{{$\varphi$}}1 {ω}{{$\omega$}}1
      % O basicstyle das listagens nao e monoespacado, entao "--" viraria uma
      % ligadura de meia-risca: `conan install --build` sairia como `-build`.
      % Esta regra mantem os dois hifens separados.
      {--}{{-{}-}}2
      % travessoes e desenho de caixa (separadores em ASCII-art)
      {—}{{---}}1 {–}{{--}}1
      {─}{{\rule[0.35ex]{0.5em}{0.4pt}}}1
      {│}{{\rule[-0.3ex]{0.4pt}{1.4ex}}}1
      {└}{{`}}1 {├}{{|}}1
      % vereditos (sim/nao/atencao) -- herdados das tabelas dos README
      {✓}{{\ding{51}}}1 {✗}{{\ding{55}}}1 {⚠}{{\textbf{!}}}1
      {↔}{{$\leftrightarrow$}}1 {∞}{{$\infty$}}1 {∈}{{$\in$}}1
      % numeros circulados usados para enumerar passos no codigo
      {①}{{\circnum{1}}}1 {②}{{\circnum{2}}}1 {③}{{\circnum{3}}}1
      {④}{{\circnum{4}}}1 {⑤}{{\circnum{5}}}1 {⑥}{{\circnum{6}}}1
      {⑦}{{\circnum{7}}}1 {⑧}{{\circnum{8}}}1 {⑨}{{\circnum{9}}}1
}

\lstdefinestyle{bt}{
    style=btchars,
    keywordstyle=\color{codekw}\bfseries,
    commentstyle=\color{codecomment}\itshape,
    stringstyle=\color{codecomment},
    breaklines=true, frame=single, framesep=4pt,
    showstringspaces=false, captionpos=b, tabsize=2,
    xleftmargin=0.075\textwidth,
    xrightmargin=0.075\textwidth
}

\tolerance=3000
\emergencystretch=3em

% =============================================================================
%  Caracteres UTF-8 fora do repertorio T1
%  O pdflatex nao conhece estes caracteres nativamente; sem os mapeamentos
%  abaixo cada ocorrencia vira um erro ("Unicode character not set up") e o
%  latexmk aborta antes de resolver sumario e referencias cruzadas.
%  Eles aparecem sobretudo em comentarios dentro de blocos lstlisting.
%  Para usar um caractere novo, acrescente uma linha aqui.
% =============================================================================
\usepackage{newunicodechar}
\usepackage{textcomp}
% Alguns destes caracteres ja tem um significado padrao valido apenas em modo
% texto; redefini-los e intencional, entao o aviso correspondente e silenciado.
\WarningFilter{newunicodechar}{Redefining Unicode character}

% Numeros circulados usados para enumerar passos dentro do codigo
\newcommand{\circnum}[1]{\raisebox{0.15ex}{\textcircled{\scriptsize #1}}}
\newunicodechar{①}{\circnum{1}}
\newunicodechar{②}{\circnum{2}}
\newunicodechar{③}{\circnum{3}}
\newunicodechar{④}{\circnum{4}}
\newunicodechar{⑤}{\circnum{5}}
\newunicodechar{⑥}{\circnum{6}}
\newunicodechar{⑦}{\circnum{7}}
\newunicodechar{⑧}{\circnum{8}}
\newunicodechar{⑨}{\circnum{9}}

% Desenho de caixa (separadores em ASCII-art dentro das listagens)
\newunicodechar{─}{\rule[0.35ex]{0.5em}{0.4pt}}
\newunicodechar{│}{\rule[-0.3ex]{0.4pt}{1.4ex}}
\newunicodechar{└}{\textnormal{`}}
\newunicodechar{├}{\textnormal{|}}

% Vereditos (sim/nao/atencao). O pifont da os dois primeiros; o terceiro nao
% tem glifo no repertorio T1 e vira uma exclamacao em negrito.
\usepackage{pifont}
\newunicodechar{✓}{\ding{51}}
\newunicodechar{✗}{\ding{55}}
\newunicodechar{⚠}{\textbf{!}}

% Simbolos matematicos e gregos
\newunicodechar{→}{\ensuremath{\rightarrow}}
\newunicodechar{↔}{\ensuremath{\leftrightarrow}}
\newunicodechar{∞}{\ensuremath{\infty}}
\newunicodechar{∈}{\ensuremath{\in}}
\newunicodechar{←}{\ensuremath{\leftarrow}}
\newunicodechar{≈}{\ensuremath{\approx}}
\newunicodechar{≤}{\ensuremath{\leq}}
\newunicodechar{≥}{\ensuremath{\geq}}
\newunicodechar{≠}{\ensuremath{\neq}}
\newunicodechar{×}{\ensuremath{\times}}
\newunicodechar{−}{\ensuremath{-}}
\newunicodechar{·}{\ensuremath{\cdot}}
% Grau e expoentes: \ensuremath para funcionarem tambem dentro de $...$
% (o significado padrao, \textdegree, e invalido em modo matematico)
\newunicodechar{°}{\ensuremath{^\circ}}
\newunicodechar{²}{\ensuremath{^2}}
\newunicodechar{³}{\ensuremath{^3}}
\newunicodechar{…}{\ensuremath{\dots}}
\newunicodechar{α}{\ensuremath{\alpha}}
\newunicodechar{β}{\ensuremath{\beta}}
\newunicodechar{Δ}{\ensuremath{\Delta}}
\newunicodechar{θ}{\ensuremath{\theta}}
\newunicodechar{λ}{\ensuremath{\lambda}}
\newunicodechar{μ}{\ensuremath{\mu}}
% U+00B5 MICRO SIGN -- visualmente identico ao mu grego, e o que os
% programas do repositorio imprimem ao escrever microssegundos.
\newunicodechar{µ}{\ensuremath{\mu}}
\newunicodechar{§}{\S}
\newunicodechar{π}{\ensuremath{\pi}}
\newunicodechar{σ}{\ensuremath{\sigma}}
\newunicodechar{τ}{\ensuremath{\tau}}
\newunicodechar{φ}{\ensuremath{\varphi}}
\newunicodechar{ω}{\ensuremath{\omega}}

% =============================================================================
%  Caixas de destaque (tcolorbox)
%
%  Quatro ambientes, cada um com um papel fixo. A regra de uso e estrita:
%
%    regra{titulo}      -- uma norma de escrita de C++ que o desenvolvedor
%                          deve seguir
%    armadilha{titulo}  -- um comportamento real da linguagem ou da biblioteca
%                          que contraria a expectativa razoavel e que o
%                          compilador nao denuncia (conversao silenciosa, move
%                          suprimido, UB, default surpreendente)
%    porque{titulo}     -- a motivacao de uma decisao de projeto
%    aviso              -- declaracao preliminar sobre a producao do documento;
%                          titulo fixo, sem argumento, de uso unico (abre a
%                          pagina 1, no topo de chapters/01-introducao.tex)
%
%  Fora o aviso unico, nada que nao se encaixe num daqueles tres papeis deve
%  virar caixa -- e `aviso` nao deve ser reaproveitado. Transcricoes
%  literais do fonte continuam sendo lstlisting; o que as identifica e a
%  legenda, que traz o caminho do arquivo dentro de src/ (convencao descrita
%  na Secao das convencoes e conferida por scripts/check-fontes.py).
% =============================================================================
\usepackage[most]{tcolorbox}

\definecolor{boxregra}{RGB}{0,39,76}       % MASA-Blue
\definecolor{boxarmadilha}{RGB}{150,40,20}
\definecolor{boxporque}{RGB}{14,90,60}
\definecolor{boxfonte}{RGB}{90,90,90}

\newtcolorbox{regra}[1][]{
    enhanced, breakable, colback=boxregra!4, colframe=boxregra,
    coltitle=white, fonttitle=\bfseries\small, title={Regra --- #1},
    boxrule=0.6pt, left=6pt, right=6pt, top=5pt, bottom=5pt,
    toprule=2pt, sharp corners=downhill}

\newtcolorbox{armadilha}[1][]{
    enhanced, breakable, colback=boxarmadilha!4, colframe=boxarmadilha,
    coltitle=white, fonttitle=\bfseries\small, title={Armadilha --- #1},
    boxrule=0.6pt, left=6pt, right=6pt, top=5pt, bottom=5pt,
    toprule=2pt, sharp corners=downhill}

\newtcolorbox{porque}[1][]{
    enhanced, breakable, colback=boxporque!4, colframe=boxporque,
    coltitle=white, fonttitle=\bfseries\small, title={Por que --- #1},
    boxrule=0.6pt, left=6pt, right=6pt, top=5pt, bottom=5pt,
    toprule=2pt, sharp corners=downhill}

\newtcolorbox{aviso}{
    enhanced, breakable, colback=boxfonte!4, colframe=boxfonte,
    coltitle=white, fonttitle=\bfseries\small, title={Aviso},
    boxrule=0.6pt, left=6pt, right=6pt, top=5pt, bottom=5pt,
    toprule=2pt, sharp corners=downhill}

% =============================================================================
%  Glossario
%
%  Usa \makenoidxglossaries (resolucao em LaTeX puro), e nao o programa externo
%  `makeglossaries`: assim o `make` continua sendo um latexmk e nada mais.
% =============================================================================
\usepackage[nopostdot,nonumberlist,toc,acronym]{glossaries}
\makenoidxglossaries
% =============================================================================
%  Glossario do relatorio
%
%  Entradas para os termos que o texto usa repetidamente e que nao sao obvios
%  para quem chega ao C++ moderno pela primeira vez. Regra de inclusao: o termo
%  aparece em mais de um capitulo E nao e auto-explicativo em portugues.
%
%  No texto, use \gls{chave} para a primeira ocorrencia relevante de um termo
%  numa secao; as demais podem ser \texttt{} normal.
% =============================================================================

% --- O modelo de objeto e de valor -------------------------------------------

\newglossaryentry{raii}{
  sort={RAII},
  name={RAII},
  description={\emph{Resource Acquisition Is Initialization}. O idioma central
    do C++: o tempo de vida de um recurso (memória, arquivo, \emph{lock},
    \emph{socket}) é amarrado ao tempo de vida de um objeto. Adquirir é
    construir; liberar é destruir. Como o destrutor roda também durante o
    desenrolamento da pilha, o idioma é o que torna exceções seguras}}

\newglossaryentry{ub}{
  sort={UB comportamento indefinido},
  name={UB (comportamento indefinido)},
  description={Situação em que o padrão não impõe nenhuma exigência ao
    programa. Não é ``dá erro'': é o compilador autorizado a supor que aquilo
    nunca acontece, e a otimizar a partir dessa suposição. Ler variável não
    inicializada, estourar um índice, violar \emph{strict aliasing} e escrever
    num objeto genuinamente \texttt{const} são casos deste relatório}}

\newglossaryentry{categoriadevalor}{
  sort={categoria de valor},
  name={categoria de valor},
  description={A classificação de uma \emph{expressão} segundo duas perguntas
    independentes --- tem identidade? pode ser movida? --- que produzem
    \emph{lvalue}, \emph{xvalue} e \emph{prvalue}. É do que depende a escolha
    entre copiar e mover, e a fonte da regra de que tudo o que tem nome é
    \emph{lvalue}}}

\newglossaryentry{move}{
  sort={move},
  name={move},
  description={Transferência do conteúdo de um objeto para outro, deixando a
    origem num estado válido porém não especificado. \texttt{std::move} não
    move nada: é um \texttt{static\_cast} para \emph{rvalue reference} que
    apenas \emph{autoriza} a operação, executada pelo construtor ou pela
    atribuição de \emph{move} do tipo}}

\newglossaryentry{elisao}{
  sort={elisao de copia},
  name={elisão de cópia},
  description={Omissão da construção de um temporário: o objeto é construído
    diretamente no destino. Desde o C++17 é \emph{obrigatória} para
    \emph{prvalues} (RVO) e opcional para variáveis locais retornadas (NRVO).
    Escrever \texttt{return std::move(local)} desliga a segunda}}

\newglossaryentry{forwarding}{
  sort={perfect forwarding},
  name={\emph{perfect forwarding}},
  description={Repasse de um argumento preservando tipo, constância e
    categoria de valor. Exige as duas peças juntas: o parâmetro \texttt{T\&\&}
    deduzido (referência universal) e o \texttt{std::forward<T>} na hora de
    repassar}}

\newglossaryentry{slicing}{
  sort={object slicing},
  name={\emph{object slicing}},
  description={Cópia de um objeto derivado para uma variável do tipo base:
    o que não cabe em \texttt{sizeof(Base)} é amputado, e com ele o
    polimorfismo. Nenhum \emph{warning} é emitido. Passar hierarquia por valor
    e usar \texttt{std::vector<Base>} são as duas formas usuais de cometê-lo}}

\newglossaryentry{vtable}{
  sort={vtable vptr},
  name={vtable / vptr},
  description={A implementação usual do despacho dinâmico: cada classe
    polimórfica tem uma tabela de ponteiros de função (\emph{vtable}) e cada
    objeto guarda um ponteiro para a tabela da sua classe real (\emph{vptr}).
    O custo é um ponteiro por objeto e uma indireção por chamada --- mas o
    preço maior costuma ser a impossibilidade de \emph{inline}}}

\newglossaryentry{rtti}{
  sort={RTTI},
  name={RTTI},
  description={\emph{Run-Time Type Information}. O mecanismo que sustenta
    \texttt{typeid} e \texttt{dynamic\_cast}. Exige pelo menos uma função
    virtual na base e pode ser desligado com \texttt{-fno-rtti}}}

% --- Genericidade e compilacao -----------------------------------------------

\newglossaryentry{sfinae}{
  sort={SFINAE},
  name={SFINAE},
  description={\emph{Substitution Failure Is Not An Error}. Se a substituição
    dos parâmetros de \emph{template} produz um tipo inválido na assinatura,
    o candidato é apenas descartado da resolução de sobrecarga em vez de gerar
    erro. Era a forma pré-C++20 de restringir \emph{templates}; \emph{concepts}
    a substitui com mensagens legíveis}}

\newglossaryentry{crtp}{
  sort={CRTP},
  name={CRTP},
  description={\emph{Curiously Recurring Template Pattern}: a classe derivada
    passa a si mesma como argumento de \emph{template} da base
    (\texttt{struct D : Base<D>}). Dá polimorfismo resolvido em compilação ---
    sem \emph{vtable}, sem indireção e com \emph{inline} --- ao preço de
    perder a base comum em tempo de execução}}

\newglossaryentry{concept}{
  sort={concept},
  name={\emph{concept}},
  description={Predicado sobre tipos, avaliado em compilação, que restringe
    quais argumentos um \emph{template} aceita. Move o erro do interior da
    instanciação para o ponto da chamada e participa da resolução de
    sobrecarga por \emph{subsumption}}}

\newglossaryentry{constexpr}{
  sort={constexpr consteval},
  name={\texttt{constexpr} / \texttt{consteval}},
  description={\texttt{constexpr} marca o que \emph{pode} ser avaliado em
    compilação; \texttt{consteval} marca o que \emph{tem} de ser. A diferença
    entre ``o compilador talvez calcule'' e ``o compilador calcula ou recusa
    o programa''}}

\newglossaryentry{typeerasure}{
  sort={type erasure},
  name={\emph{type erasure}},
  description={Guardar objetos de tipos não relacionados atrás de uma
    interface uniforme, apagando o tipo estático mas preservando o
    comportamento. \texttt{std::function}, \texttt{std::any} e o
    \texttt{std::type\_index} do ECS deste repositório são três realizações
    diferentes da mesma ideia}}

% --- Biblioteca padrao -------------------------------------------------------

\newglossaryentry{iterador}{
  sort={iterador},
  name={iterador},
  description={A abstração que separa algoritmos de containers: um objeto que
    aponta para um elemento e sabe avançar. As categorias (\emph{input},
    \emph{forward}, \emph{bidirectional}, \emph{random access}) definem que
    algoritmos podem operar sobre qual container e a que custo}}

\newglossaryentry{view}{
  sort={view},
  name={\emph{view} (C++20)},
  description={Um \emph{range} que não possui os elementos e cujas operações
    de cópia e destruição são O(1). \emph{Views} são preguiçosas: compor
    \texttt{filter} com \texttt{transform} não percorre nada até que alguém
    itere o resultado}}

\newglossaryentry{invalidacao}{
  sort={invalidacao de iteradores},
  name={invalidação de iteradores},
  description={A perda de validade de iteradores, ponteiros e referências
    após uma modificação do container. As regras variam por container e são a
    origem mais comum de \gls{ub} em código que ``funcionava''}}

\newglossaryentry{pmr}{
  sort={pmr},
  name={\texttt{std::pmr}},
  description={\emph{Polymorphic Memory Resources} (C++17). Permite trocar a
    estratégia de alocação de um container em tempo de execução, sem que a
    escolha entre no tipo --- ao contrário do parâmetro \texttt{Allocator}
    clássico, que faz \texttt{vector<int,A>} e \texttt{vector<int,B>} serem
    tipos distintos}}

% --- Concorrencia ------------------------------------------------------------

\newglossaryentry{datarace}{
  sort={data race},
  name={\emph{data race}},
  description={Dois acessos concorrentes ao mesmo endereço, ao menos um deles
    de escrita, sem sincronização entre eles. É \gls{ub} --- não ``valor
    errado ocasional''. O detector prático é o ThreadSanitizer}}

\newglossaryentry{atomico}{
  sort={operacao atomica},
  name={operação atômica},
  description={Operação indivisível sobre uma variável compartilhada, com uma
    ordem de memória associada. Resolve a corrida sobre \emph{aquela}
    variável; não estabelece invariantes entre duas variáveis --- para isso
    ainda é preciso um \emph{mutex}}}

\newglossaryentry{corotina}{
  sort={co-rotina},
  name={co-rotina (C++20)},
  description={Função que pode suspender a própria execução e ser retomada
    depois, preservando o estado local num \emph{frame} alocado à parte.
    A linguagem fornece só o mecanismo (\texttt{co\_await},
    \texttt{co\_yield}, \texttt{co\_return}); os tipos \texttt{generator} e
    \texttt{task} são escritos pelo usuário ou pela biblioteca}}

% --- Ferramentas e ecossistema -----------------------------------------------

\newglossaryentry{meson}{
  sort={Meson},
  name={Meson},
  description={O sistema de \emph{build} deste repositório. Descreve o projeto
    numa DSL declarativa e gera arquivos do Ninja; a compilação em si é do
    Ninja. Todo subprojeto tem um \texttt{meson.build} próprio, registrado em
    \texttt{src/meson.build}}}

\newglossaryentry{conan}{
  sort={Conan},
  name={Conan 2},
  description={O gerenciador de dependências do repositório. Lê o
    \texttt{conanfile.py}, resolve e compila as bibliotecas de terceiros e
    gera os arquivos que fazem o \texttt{dependency()} do Meson encontrá-las}}

\newglossaryentry{abi}{
  sort={ABI},
  name={ABI},
  description={\emph{Application Binary Interface}: o contrato em nível de
    binário --- convenção de chamada, \emph{layout} de objetos, decoração de
    nomes. É o que decide se dois \texttt{.o} compilados separadamente podem
    ser ligados, e o que \texttt{extern "C"} e o ABI System V x86-64 fixam}}


% =============================================================================
%  Indice remissivo de classes e portas
%
%  \classe{std::vector}     imprime \texttt{std::vector} e indexa
%  \cabecalho{memory}       imprime \texttt{<memory>} e indexa sob "cabecalhos"
%  \projeto{value_semantics} imprime \texttt{src/value\_semantics/} e indexa
%                            sob "subprojetos"
%  Os tres sao apenas atalhos: nada impede escrever \texttt{} direto quando a
%  ocorrencia nao merece entrada no indice.
% =============================================================================
% Regras horizontais de qualidade tipografica para as tabelas de referencia
\usepackage{booktabs}

\usepackage{makeidx}
\makeindex

\newcommand{\classe}[1]{\texttt{#1}\index{#1@\texttt{#1}}}
\newcommand{\cabecalho}[1]{\texttt{<#1>}\index{cabeçalhos!#1@\texttt{<#1>}}}
\newcommand{\projeto}[1]{\texttt{src/#1/}\index{subprojetos!#1@\texttt{#1}}}

% =============================================================================
%  Diagramas em TikZ
%
%  Todas as figuras deste manual sao TikZ escrito a mao dentro do proprio
%  capitulo: nao ha PNG, nao ha pipeline de renderizacao externa e nao ha
%  dependencia de navegador para regerar uma figura. Editar o desenho e editar
%  o .tex ao lado do texto que ele ilustra.
% =============================================================================
\usepackage{tikz}
\usetikzlibrary{arrows.meta,positioning,fit,backgrounds,calc,shapes.geometric,
                shapes.multipart,decorations.pathreplacing,matrix}

% Estilos comuns aos diagramas do manual.
\tikzset{
    caixa/.style     = {draw=boxregra, fill=boxregra!6, rounded corners=2pt,
                        align=center, inner sep=4pt, font=\footnotesize},
    caixaforte/.style= {draw=boxregra, fill=boxregra!18, rounded corners=2pt,
                        align=center, inner sep=4pt, font=\footnotesize\bfseries},
    caixafraca/.style= {draw=gray!60, fill=gray!8, rounded corners=2pt,
                        align=center, inner sep=4pt, font=\footnotesize},
    caixaok/.style   = {draw=boxporque, fill=boxporque!8, rounded corners=2pt,
                        align=center, inner sep=4pt, font=\footnotesize},
    caixarisco/.style= {draw=boxarmadilha, fill=boxarmadilha!8, rounded corners=2pt,
                        align=center, inner sep=4pt, font=\footnotesize},
    folha/.style     = {draw=boxregra, fill=white, rounded corners=6pt,
                        align=center, inner sep=4pt, font=\footnotesize},
    seta/.style      = {-{Stealth[length=2mm]}, draw=boxregra!80},
    setafraca/.style = {-{Stealth[length=2mm]}, draw=gray!70, dashed},
    setaok/.style    = {-{Stealth[length=2mm]}, draw=boxporque!85},
    setarisco/.style = {-{Stealth[length=2mm]}, draw=boxarmadilha!85},
    rotulo/.style    = {font=\scriptsize\itshape, inner sep=1pt},
    aresta/.style    = {draw=boxregra!70},
}

% =============================================================================
%  hyperref -- SEMPRE por ultimo (redefine comandos de referencia).
%  Da sumario clicavel e painel de marcadores no visualizador, o que torna a
%  pre-visualizacao local navegavel. `hidelinks` mantem a aparencia impressa
%  (sem molduras nem cores); troque por colorlinks=true se preferir links
%  coloridos na tela.
% =============================================================================
\usepackage{hyperref}
\hypersetup{
    hidelinks,
    bookmarksnumbered = true,
    bookmarksopen     = true,
    pdftitle          = {\mytitle},
    pdfauthor         = {\myauthors},
}

% cleveref DEPOIS de hyperref: referencias que ja trazem o nome do objeto,
% de modo que \cref{sec:modelo-de-nos} imprime "Secao III" sem repeticao
% manual do rotulo em cada chamada.
% Sem opcao de idioma: o documento nao carrega babel, entao os nomes vem das
% declaracoes \crefname abaixo.
\usepackage{cleveref}

% O documento nao carrega babel, entao os rotulos automaticos vem em ingles
% ("Figure 5", "Table 2", "Contents"). Traduzi-los na mao e a unica opcao --
% carregar babel[brazil] aqui quebraria o template MASA, que fixa fontes e
% espacamentos supondo o padrao.
\renewcommand{\figurename}{Figura}
\renewcommand{\tablename}{Tabela}
\renewcommand{\contentsname}{Sumário}
\renewcommand{\listfigurename}{Lista de figuras}
\renewcommand{\listtablename}{Lista de tabelas}
\renewcommand{\lstlistingname}{Listagem}
\renewcommand{\lstlistlistingname}{Lista de listagens}

\crefname{section}{Seção}{Seções}
\Crefname{section}{Seção}{Seções}
\crefname{figure}{Figura}{Figuras}
\Crefname{figure}{Figura}{Figuras}
\crefname{table}{Tabela}{Tabelas}
\Crefname{table}{Tabela}{Tabelas}
\crefname{equation}{Equação}{Equações}
\Crefname{equation}{Equação}{Equações}
\crefname{lstlisting}{Listagem}{Listagens}
\Crefname{lstlisting}{Listagem}{Listagens}

% Diretorio de busca de imagens. O manual nao usa \includegraphics no corpo
% (as figuras sao TikZ), mas a capa do template MASA carrega os PDF de
% assets/, entao o caminho precisa continuar resolvendo.
\graphicspath{%
  {figures/}%
  {assets/}%
}
